\documentclass[a4paper, 12pt, oneside, BCOR=1cm,toc=chapterentrywithdots]{scrbook}

\usepackage{graphicx}                   % use for pdfLatex
\usepackage{makeidx}                    % \printindex
\usepackage[colorlinks=false]{hyperref}
% \usepackage{tocbibind}
\usepackage{blindtext}
\usepackage{subfigure}
\usepackage{acronym}
\usepackage[german]{babel}
\usepackage{csquotes}
\usepackage[backend=biber,style=alphabetic]{biblatex}
\addbibresource{bibliography.bib}

\defbibfilter{papers}{
  type=article or
  type=inproceedings
}

\hypersetup{bookmarksnumbered=true,bookmarksopen=true, bookmarksopenlevel=1,pdfborder=0 0 0}

% remove this later
\setcounter{tocdepth}{4}

\begin{document}
\begin{titlepage}

  \begin{center}
    \raisebox{-1ex}{\includegraphics[scale=1.5]{assets/tuc.pdf}}\\
  \end{center}
  \vspace{0.5cm}

  \begin{center}

    \LARGE{\textbf{Title of the bachelor thesis}}\\
    \vspace{1cm}


    \Large{\textbf{Bachelor Thesis}}\\
    \vspace{1cm}
    Submitted in Fulfilment of the\\
    Requirements for the Academic Degree\\
    B.Sc.\\
    \vspace{0.5cm}
    Dept. of Computer Science\\
    Chair of Computer Engineering
  \end{center}
  \vspace{3cm}
  Eingereicht von: Joshua Jeschek\\
  Matrikelnummer: 652272\\
  Datum: 24.12.2023\\
  \vspace{0.3cm}\\
  Dr. Thomas Wilhelm-Stein \\
  Dr. Heino Rudolf

\end{titlepage}

\addchap*{Abstract}
\blindtext


\tableofcontents

\begingroup
\cleardoublepage
\addcontentsline{toc}{chapter}{\listfigurename}
\listoffigures

\let\clearpage\relax

\cleardoublepage
\addcontentsline{toc}{chapter}{\listtablename}
\listoftables
\endgroup

\twocolumn
% \addchap{List of Abbreviations}
% \begin{acronym}[Bash]
% 	\acro{SQL}{Structured Query Language}
% \end{acronym}

\onecolumn

\chapter{Einleitung}

\section{Hintergrund und Motivation}
\section{Ziele und Umfang}
\section{Aufbau der Arbeit}

\chapter{Theoretische Grundlagen}

\section{envVisio}
\subsection{Das Realitätsmodell}
\subsubsection{Allgemeine Erklärung von envVisio}
\subsubsection{Vielleicht Aufteilung in mehrere Abschnitte?}
\subsubsection{was ist alles wichtig?}
\subsection{Laden und Konvertieren}
\subsubsection{auf Loader und insbesondere Converter eingehen}
\subsubsection{da in Converter Formular eingesetzt werden soll}
\subsubsection{wichtig, Prozesse zu verstehen}

\section{Human-Centered Design}
\subsubsection{Begriffe und Konzepte aus Design of Everyday Things}
\subsubsection{Abschnittstitel nicht final, vielleicht is HCD auch ein Unterpunkt auf gleicher Höhe wie Aff., Sign. und Map.}
\subsection{Affordances, Signifiers und Mappings}
\subsubsection{Kann gut bei 2.3.2 angewendet werden}
\subsection{Activity-Centered Design}

\section{Formulare}
\cite{wroblewskiWebForm2008}
\subsection{Funktionsweise}
\subsubsection{Technische Fkt.-weise von F.}
\subsection{Designanfordungen}
\cite{normanDesignEveryday2013} (auf Affordances, Signifiers und Mappings eingehen)
\subsubsection{Rückbezug auf 2.2}
\subsubsection{Anordnung}
\subsubsection{Kontrollelemente}
\subsubsection{Bezeichnungen und Platzhalter}
\subsubsection{Validierung}

\section{Usability Testing}
\cite{barnumUsabilityTesting2021}
\subsection{Testmethoden}
\subsubsection{Quantitative vs Qualitative Tests}
\subsubsection{Formative vs Summative Studien}
\subsection{Teilnehmer:innen}
\subsubsection{Wer wird ausgewählt?}
\subsection{Fragebogen}
\subsubsection{Wie können Fragebögen zur Evaluierung erstellt werden?}
\subsubsection{Verwendung von Standard-Fragebogen}

\section{Analyse und Vergleich von existierenden Projekten}
\subsection{Home Assistant}
\subsection{DBSnap}
\url{https://www.semanticscholar.org/paper/DBSnap\%3A-Learning-Database-Queries-by-Snapping-Silva-Chon/aabf7791575f8cc62432ca5e6289862eec0f6845}
\url{https://www.semanticscholar.org/paper/DBSnap-2\%3A-New-Features-to-Construct-Database-by-Silva-Loza/61c95757d7d4151f51ea848623c6839816756385}
\subsection{SQheLper}
\url{https://www.semanticscholar.org/paper/SQheLper\%3A-A-block-based-syntax-support-for-SQL-Jacobs-Jaschke/b4f6023a4ee72d0103ee521c4f9b35768d38810b}
\subsection{Scratch}


\chapter{Implementierung}
\section{Zielformulierung}
\subsubsection{Genaue Formulierung der Problemstellung, was ist das Ziel?}
\section{Lösungsvorstellung}
\subsubsection{Erklärung des Resultats}

\chapter{Studie}
\section{Studienaufbau}
\subsection{Auswahl von Teilnehmer:innen}
\subsection{Fragebogen}
\subsection{Studienablauf}
\section{Ergebnisse und Analyse}
\section{Schlussfolgerungen und Zukunftsaussicht}

\chapter{Fazit}


\nocite{*}
\chapter*{Quellen}
\begingroup
\let\clearpage\relax
\printbibliography[type=article, title=Arbeiten?, heading=subbibliography]
\printbibliography[type=book, title=Bücher, heading=subbibliography]
\printbibliography[type=misc, title=sonstige, heading=subbibliography]
\endgroup

\printindex

\end{document}
